\section{Discussion}
\label{sec:conclusion}

The clean finding is that, in our tree-regularized prototype model
on a real hierarchical-classification problem, the latent manifold
has a local effect on tree fidelity even when its global effect is
unstable. Hyperbolic prototype geometry is the only one of
our trained models that meaningfully improves on the frozen encoder
for local hierarchy preservation. Sibling recall@5 at $d{=}8$ is
$0.195$ for hyperbolic, $0.149$ for Euclidean, and $0.160$ for
$k$-NN-5 on raw CLIP features: the Euclidean fit is tied with the
encoder; the hyperbolic fit improves on it by
$3.5$\,pp. The advantage holds when sibling sets are redefined from
data-driven empirical trees built from CLIP or DINOv2 features, and
aggregates to $+8.7$\,pp on sibling and $+15.2$\,pp on cousin recall
over the full dimension sweep with paired-$t$ $p<10^{-4}$.

The classification gap is real but uninformative about geometry. A
logistic-regression head on raw CLIP features achieves $64.1\%$
top-1, which the Euclidean fit matches at $d{=}8$ ($64.3\%$) and
the hyperbolic fit underperforms by $4$--$6$\,pp. Neither fit is doing meaningful
classification work that a linear head on the same features cannot;
the difference between the two fits on this axis is best read as
``Euclidean prototype distances do not distort classification beyond
what the encoder already supports, while hyperbolic prototype
distances do.''

The global tree-fidelity comparison is the place where careful
framing matters most. Mean tree distortion favours Euclidean
significantly against the default tree; class-center / tree Spearman
is not significantly different from zero across the sweep
($p=0.68$); the small global-Spearman gap reverses sign between
empirical reference trees built from CLIP and from DINOv2 feature
centroids. ``Either geometry wins on global tree fidelity'' is not a
claim our experiments support. Sibling and cousin recall favour
hyperbolic on every reference tree we tried.

\paragraph{Limitations.} We cannot generalise beyond medium-scale,
Western-canon-heavy WikiArt with a frozen CLIP-ViT-B/16 encoder for
training; the local advantage we identify may not survive a
fine-tuned encoder or a hierarchy substantially deeper than the
$27$-leaf taxonomy used here. We cannot conclude that the hyperbolic
geometry does substantive work inside the head MLP, since only the
prototypes live on the manifold.

\paragraph{Extensions.} Three natural follow-ups in increasing cost.
A genuinely hyperbolic head along the lines of M\"obius layers
\citep{ganea2018hyperbolic} would let the manifold shape the
representation rather than only the decision boundary, and is the
cheapest test of whether the classification gap closes. The
empirical reference trees in this paper are built in Euclidean space
from encoder features. Building one in hyperbolic space, or by direct
tree-learning methods \citep{chami2020trees}, would close a
remaining circularity. Replication on a deeper hierarchy
(iNaturalist with WordNet, or a fine-grained subdivision of
WikiArt's largest classes) would test the limits of the local
advantage we observe.
